\section{Computation and source comparison}
\subsection{Exact arithmetic and independent checks}
The accompanying package separates three kinds of evidence:
\begin{enumerate}[leftmargin=*,itemsep=3pt]
 \item Integer and rational identities are checked with Python integers and
 \texttt{Fraction}. These include exact counts, defect polynomials,
 Poisson transforms, marked coefficients, and critical coefficients.
 \item Transcendental evaluations use optional \texttt{mpmath} arithmetic.
 They corroborate the theorems and expose finite-$n$ limitations; they
 are not certified interval enclosures.
 \item The proofs above provide asymptotic remainders and quantifiers.
 Neither fitted coefficients nor numerical agreement replaces those proofs.
\end{enumerate}
The exact core requires no symbolic algebra package. It computes finite
power sums by telescoping, the $H_j$ recurrence by rational polynomial
arithmetic, and Poisson transforms by the Touchard recurrence.
The critical generator separately performs weighted polynomial truncation
and exact centered-moment substitution.

Independent checks reconstruct polynomial expectations in the
falling-factorial basis. If a polynomial $H$ has degree at most $d$, then
\begin{equation}\label{spb:ep:eq:diffbasis}
 H(r)=\sum_{j=0}^d\Delta^jH(0)\binom rj,
 \qquad \E H(X)=\sum_{j=0}^d\Delta^jH(0)\frac{\lambda^j}{j!},
\end{equation}
where $\Delta H(r)=H(r+1)-H(r)$.
This is a distinct route from converting ordinary powers through Touchard
polynomials. Extra integer samples check the interpolated identities.
The public tests also multiply the formal generating function independently
at small sizes, check every residue for $p=1,2,3,4$ in numerical identity
tests, and test active input guards under optimized Python.

Program limits are deliberately finite resource safeguards, not restrictions
of the theorems. For example, the public generator caps the formal $H$ order
at 16, fractional order at 32 subject to that cap, and critical order at 8.
All guards use ordinary conditionals and exceptions, so Python's
\texttt{-O} flag cannot remove them. The precise limits and commands are
listed in the code README.

\subsection{Complete sums and a visible preasymptotic warning}
\label{spb:ep:sec:numerics}
For the complete-sum diagnostic, the original expression
\eqref{spb:ep:eq:counts} is evaluated as an exact integer. Its normalization uses
90-digit arithmetic. Independently, the transformed expression sums
\emph{every} allowed $r$ from zero through $n$, using logarithms and
log-gamma evaluations. No Poisson window and no asymptotic coefficients
enter either finite-sum comparison.

For $p\ge2$ let
\[
 A_K(n)=\sum_{k=0}^K C_{p,k}(c_p)n^{-k/q},\qquad
 S_J(n)=\sum_{j=0}^J U_j(\lambda)n^{-j}.
\]
These are different finite approximations. In particular, $S_{12}$ retains
many fractional powers beyond those in $A_{12}$.
Table~\ref{spb:ep:tab:counts} reports relative errors as approximation divided by
$\Rcal_p(n)$ minus one. The number of printed digits is deliberately smaller
than the working precision.

\begin{table}[htbp]
\centering\small
\begin{tabular}{@{}crrrr@{}}
\toprule
$p$ & $n$ & $\Rcal_p(n)$ & $A_{12}/\Rcal_p-1$ & $S_{12}/\Rcal_p-1$\\
\midrule
2 & 10000 & $0.5075173547743$ & $-9.64017\cdot10^{-9}$ & $+6.40647\cdot10^{-10}$\\
3 & 10000 & $0.1416076329487$ & $-0.6615364472$ & $+0.0012291163$\\
\bottomrule
\end{tabular}
\caption{Complete-sum diagnostics. Both parameterizations are asymptotically
valid at fixed order, but have sharply different errors here.}
\label{spb:ep:tab:counts}
\end{table}

The original and transformed evaluations agree to about 86 decimal places
in these cases. Such agreement checks the exact transformation numerically;
it does not certify the final transcendental digits by outward rounding.
Most importantly, the $p=3$ fractional truncation is still wrong by about
$66\%$ at $n=10000$. The shorter $t^2$ and $t^4$ approximations are worse.
Large constants such as $e^{9/2}$ delay the useful asymptotic regime.
The smaller error of $S_{12}$ at this input does not establish a general
optimal truncation principle.

For the critical case, the complete count at $n=10000$ gives
\[
 a_1(n)/M_1(n)\approx1.02446126439771172897.
\]
The half-power truncation through $C_{1,4}$ has relative error
approximately $-2.32029\cdot10^{-9}$.
The companion odd input $n=10001$ gives an error of the same size,
consistent with the common algebraic coefficients for both parities.
\begin{writenote}[The decimals recomputed, 7 October 2026]
The write recomputed Table~\ref{spb:ep:tab:counts} and the critical values
from the exact integers at 110 digits, with its own coefficients
$C_{p,k}$, $k\le12$, and transforms $U_j$, $j\le12$: every printed value is
the correct rounding of the recomputed one. Several are roundings rather than
truncations; the truncations at the printed length are $\Rcal_2=0.5075173547742\ldots$,
$A_{12}/\Rcal_2-1=-9.64016\ldots\cdot10^{-9}$,
$S_{12}/\Rcal_2-1=+6.40646\ldots\cdot10^{-10}$,
$\Rcal_3=0.1416076329486\ldots$, $A_{12}/\Rcal_3-1=-0.6615364471\ldots$,
$S_{12}/\Rcal_3-1=+0.0012291162\ldots$, $a_1(n)/M_1(n)=1.02446126439771172896\ldots$
and the relative error $-2.32028\ldots\cdot10^{-9}$ at $n=10000$
($-2.31971\ldots\cdot10^{-9}$ at $n=10001$). The integer and transformed
evaluations agree to $10^{-106}$ in the write's computation.
\end{writenote}

\subsection{Marked and inverse diagnostics}
At $n=10^9$ the reported TV distances divided by their proved leading
equivalents are approximately
\begin{center}
\begin{tabular}{@{}crr@{}}
\toprule
$p$ & Original conditioned reference & Tilted conditioned reference\\
\midrule
2 & $1.000111734$ & $1.000214479$\\
3 & $1.000578191$ & $0.999189636$\\
\bottomrule
\end{tabular}
\end{center}
These computations use central windows, with the window endpoints and
number of lattice points included in the machine-readable output.
Chernoff bounds control the difference between the window-normalized TV
and complete-law TV. Their high-precision evaluated bounds are tiny, but
are still not interval certificates. The moment outputs are explicitly
window diagnostics; a TV bound alone does not control unbounded moment
errors.

The Newton diagnostic has two modes. By default it sets $Y=F_K(n)$, so the
finite model root is known exactly to be $n$; this isolates iteration error.
An optional exact-target mode instead uses $Y=\log a_p(n)$ and therefore
includes model truncation error. The program makes exactly the requested
finite number of Newton steps and checks numerical domain conditions at
each step. Neither mode supplies a certified convexity onset or certifies
threshold rounding. These distinctions matter when interpreting a tiny
printed residual.
\begin{writenote}[The marked ratios recomputed, 7 October 2026]
With the exact law over a full central window (log-gamma values of the
summands) and the residue-conditioned Poisson references, the write obtains
the four ratios at $n=10^9$ as $1.00011173445\ldots$, $1.00021447898\ldots$
($p=2$) and $1.00057819116\ldots$, $0.999189635789\ldots$ ($p=3$): the printed
values are correct roundings (truncations $1.000214478$ and $0.999189635$ for
the two that differ).
\end{writenote}

\subsection{Observed primary formulas}
\label{spb:ep:sec:sources}
The exact finite sums and the three displayed OEIS expressions were
observed in primary-entry text on 5 October 2026.
The A360479 ordinary page was retrieved directly; A360592 and A360747
were obtained through indexed primary OEIS text when direct opening was
unreliable. Independent retrievals corroborated the mathematical formulas.
No verified sequence revision identifier or complete history export was
obtained. Global page footers and indexed modification dates do not
establish the current sequence revision.

A360592 identifies V\'aclav Kote\v{s}ovec as the sequence author, dated
13 February 2023. A360479 and A360747 identify Seiichi Manyama as sequence
author, dated 19 February 2023; their refinement formulas are attributed
to Kote\v{s}ovec on 19 and 20 February 2023 respectively.
These are observed source attributions, not a reconstructed publication
history. The inspected entry links were numerical tables or graphs rather
than proof articles.

Using the normalizations of this article, the observed critical multiplier
was
\[
 1+\left(\frac c4+\frac{11c^3}{8}+\frac1{8c}\right)t,
 \qquad c=\sqrt{e/2},\quad t=n^{-1/2}.
\]
For $p=2$, it was
\begin{align*}
 1&+\left(-\frac{13c^2}{6}+\frac1{8c}\right)t\\
  &+\left(\frac{169c^4}{72}-\frac{15c}{16}
                               +\frac{67}{384c^2}\right)t^2\\
  &+\left(-\frac{2197c^6}{1296}+\frac{7913c^3}{576}
                +\frac{3929}{2304}+\frac{497}{3072c^3}\right)t^3,
\end{align*}
with $c=e^{4/3}/3^{1/3}$ and $t=n^{-1/3}$.
For $p=3$, it was
\begin{align*}
 1&+\frac{t}{8c}
 +\left(-\frac{37c^2}{8}+\frac{67}{384c^2}\right)t^2\\
  &+\left(-\frac{205c}{64}+\frac{497}{3072c^3}\right)t^3\\
  &+\left(\frac{1369c^4}{128}+\frac{10721}{3072}
                      +\frac{218831}{1474560c^4}\right)t^4,
\end{align*}
with $c=e^{9/4}/\sqrt2$ and $t=n^{-1/4}$.
Each appeared as part of a formula using $\sim$.
Every one of these multipliers tends to one, so the formulas remain
correct if read solely as leading equivalents. Their intended refined
coefficients differ from \eqref{spb:ep:eq:critC1}, \eqref{spb:ep:eq:p2}, and
\eqref{spb:ep:eq:p3}. Appendix~\ref{spb:ep:app:discrepancies} lists the differences explicitly.
No explanation of how the discrepancies arose is assumed.
\begin{writenote}[7 October 2026]
The write compared these three multipliers with the formulas of the current
entries (revisions in Remark~\ref{spb:ep:rem:oeis}): after the normalizations
of this article they are exactly the displayed expressions (checked
symbolically, leading factors included).
\end{writenote}

\subsection{Bounded historical scope}
The source comparison establishes a specific coefficient audit, not a
universal priority result. Exact-number and formula-oriented searches did
not locate a proof of the observed refinements. Two inspected default-branch
catalogues in the public ProveIt repository and exact code searches did
not identify this endpoint family. Those are bounded negative checks;
they cannot rule out an unnamed general theorem or material elsewhere.

A potentially relevant 2013 note by Kote\v{s}ovec,
\emph{Interesting asymptotic formulas for binomial sums}, was bibliographically
identified at
\url{https://www.kotesovec.cz/math_articles/kotesovec_interesting_asymptotic_formulas.pdf},
but retrieval timed out and its theorem content was not compared.
Likewise the latest OEIS histories were not verified.
Consequently no assertion that all prior relevant work has been excluded
is made here. The correction follows from the exact finite sums and the
self-contained proofs, regardless of broader historical priority.

For classical background, the DLMF inverse-function and Charlier generating
function entries were checked directly. The Lambert-$W$ reference was
verified bibliographically and at its definition-level abstract; we do
not attribute an uninspected theorem in that paper to the present family.
The analytic inverse estimates needed here are proved in
Section~\ref{spb:ep:sec:inverse}.

\begin{remark}[{[write]} The OEIS entries, 7~October 2026]\label{spb:ep:rem:oeis}
The write read the three entries in the internal format on 7~October 2026.
\emph{A360592}, revision \#28 (17 February 2023), by V\'aclav Kote\v{s}ovec
(13 February 2023), ``G.f.: Sum\_\{k>=0\} (1 + k*x)\^{}k * x\^{}k.'', with the
finite sum, a b-file by Kote\v{s}ovec for $n=0,\ldots,760$, and the unsigned
formula line $a(n)\sim\exp(\exp(1/2)\sqrt{n/2}-3e/8)\,n^{n/2}/2^{n/2+1}\,
(1+((e^{1/2}+e^{-1/2})/2^{5/2}+11e^{3/2}/2^{9/2})/\sqrt n)$.
\emph{A360479}, revision \#40 (19 February 2023), by Seiichi Manyama
(19 February 2023), with a b-file by Kote\v{s}ovec ($n=0,\ldots,635$), a graph
of the asymptotic ratio, and the refinement through $n^{-1}$ signed by
Kote\v{s}ovec, 19 February 2023. \emph{A360747}, revision \#22 (20 February
2023), by Manyama (19 February 2023), with a b-file by Kote\v{s}ovec
($n=0,\ldots,600$), a graph, and the refinement through $n^{-1}$ signed by
Kote\v{s}ovec, 20 February 2023, ending ``see graph for more minor asymptotic
terms''. These are the latest revisions, all from February 2023, so the
formulas the source compared in Section~\ref{spb:ep:sec:sources} are those of
the current entries, and its attributions are as stated there.
\emph{Checks.} All terms of the three b-files equal the write's
evaluation of~\eqref{spb:ep:eq:counts}. Rewritten in this article's
normalizations, the three formula lines are exactly the multipliers displayed
in Section~\ref{spb:ep:sec:sources}, with leading factors $M_1(n)$, $M_2(n)$,
$M_3(n)$ (SymPy). The leading equivalents are therefore correct. The first
corrections are wrong: by Theorems~\ref{spb:ep:thm:critical}
and~\ref{spb:ep:thm:allp} the true first coefficients are $C_{1,1}$,
$C_{2,1}=-13c^2/6$ and $C_{3,1}=0$, and each entry adds $1/(8c_p)$. Against
the exact sums the write finds $(\text{entry multiplier}-a_p(n)/M_p(n))/t$
equal to $0.10722073\ldots$ at $n=10^{14}$ ($p=1$; $1/(8c)=0.10722048\ldots$),
$0.04752092\ldots$ at $n=10^{18}$ ($p=2$; $1/(8c_2)=0.04752160\ldots$) and
$0.01863213\ldots$ at $n=10^{24}$ ($p=3$; $1/(8c_3)=0.01863212\ldots$), the
full refined multipliers being used; so the entries' first corrections are
refuted by the proofs of this Part and, numerically, by the exact sums. No OEIS
edit is part of this work.
\end{remark}
